The JCRIN / PCS Epoch Table as the Discrete Chronological Scaffold of the Decompositional Hierarchy

In the architecture of Joshua Christopher Ryan’s Infinitesimal Numberline (JCRIN) and its Precision Control Sequence (PCS), cosmology is no longer treated as a continuous fluid evolving under differential equations alone. It is reconstructed as a countable lattice of \(10^9\) microscopic ticks, each of length \(\varepsilon = 10^{-9}\). Within this lattice the nine incremental metadata ranges—R1 through R9—function as the primary chronological scaffold. They are not arbitrary partitions; they are the deliberate quantization of cosmic history into precision-controlled epochs that simultaneously govern numerical stability, angular resolution, and the relevance of tensor modes to the observable B-mode sky.

The table begins at the high-redshift frontier. Range R1 occupies the interval \(n = 1\,000\,000\) to \(1\,999\,999\), corresponding to redshifts \(z \approx 999 \to 500.5\). Here the universe is still emerging from the dark ages; the era is labeled “Early matter” and the scientific relevance is designated “Reionization precursor.” Six decimal places of floating-point precision are retained because the dynamical range is extreme and the risk of truncation error is acute. In the language of the discrete transport equation the complex vacuum state \(z_n = y_n - i\pi y_n\) is still far from the phase-lock horizon; the imaginary component has not yet accumulated the half-twist that will later regularize the Schwarzschild pathway. R1 therefore serves as the numerical and conceptual foundation upon which all subsequent structure is erected.

As the lattice index advances, the next three ranges—R2, R3 and R4—collapse into a single narrative unit. Spanning \(n = 2\,000\,000\) to \(4\,999\,999\) and redshifts \(z \approx 499 \to 200\), they occupy the heart of matter domination. The module itself prints them collectively as “Tensor mode growth.” Six decimal places remain mandatory; the growth of primordial gravitational-wave amplitude is still linear and the angular contribution of each range is fixed at \(\Delta\theta = 0.036^\circ\). These three ranges are the engine room of the tensor-to-scalar ratio. Their cumulative footprint begins to build the total sky area \(\Delta\Theta_{\rm total} = 0.324^\circ\) that will later be rescaled to the fiducial value \(r_{\rm fid} = 0.01\). In the broader Decompositional Hierarchy they correspond to the epoch in which the bosonic and fermionic embeddings of the complex pathway remain phase-coherent with the expanding scale factor \(a(N) = e^{\varepsilon N}\).

The final five ranges—R5 through R9—mark the transition into the late-matter universe. From \(n = 5\,000\,000\) to \(9\,999\,999\) and redshifts \(z \approx 199 \to 100\), the table declares “Peak B-mode signal.” Precision is deliberately relaxed: six decimals give way to four and then to three. This is not a loss of rigor but an adaptive quantization. As the dynamical range contracts, fewer significant figures are required to maintain machine-precision closure of the global volume constraint \(C \to 1\). The same \(\Delta\theta = 0.036^\circ\) continues to accumulate, yet the scientific emphasis shifts from the generation of tensor modes to their observational imprint on the polarized microwave sky. These ranges sit just above the redshift window where the late-time geometric torque of the Cosmological Clock will later activate; the suppression functions \(S(\chi)\), \(f_{\rm damp}(z)\) and \(f_{\rm mod}(z)\) still hold the torque in abeyance, preserving the early-universe calibration.

Taken together, the nine ranges constitute a single continuous narrative arc. They begin in the reionization precursor, pass through the linear growth of tensor modes, and culminate at the threshold of the peak B-mode window—all while remaining strictly internal to the discrete lattice that terminates at \(N = 10^9\). Every entry carries four interlocking attributes: the integer bounds on \(n\), the floating-point decimal quota required for numerical stability, the corresponding redshift interval, and the physical relevance to gravitational-wave cosmology. The angular calibration constants bind the table to the photometric interface: nine identical contributions of \(0.036^\circ\) produce a total footprint of \(0.324^\circ\), which is then rescaled by the map scalar factor so that the aggregate tensor amplitude matches the fiducial \(r = 0.01\).

In the peer-reviewed sense the table performs three indispensable functions. First, it supplies an explicit, reproducible quantization scheme that eliminates uncontrolled floating-point drift across thirteen orders of magnitude in redshift. Second, it embeds the generation and observational relevance of primordial B-modes inside the same discrete transport geometry that regularizes the Schwarzschild pathway and generates the late-time Hubble torque. Third, it renders the entire construction falsifiable: any alteration of the decimal schedule, the angular increment, or the redshift boundaries must propagate consistently through the complex vacuum states, the phase-lock limit \(\arctan(2\pi)\), and the photometric distances \(\chi(z)\) and \(d_L(z)\). The binds that hold the narrative together—modular precision, emergent residual geometry, closed algebraic identities, and directed cross-reference to the Cosmological Clock—are therefore the same binds that support rigorous external examination.

Thus the JCRIN / PCS Epoch Table is not an auxiliary lookup list. It is the chronological skeleton of the Decompositional Hierarchy: a nine-stage lattice pathway that carries the universe from the reionization precursor to the threshold of the B-mode peak while preserving, at every step, the numerical exactitude and topological coherence demanded by a self-consistent discrete cosmology.
